\SourceLeaf{166}
Unterschied zwischen den jetzigen Witwen und den künftig hinzukommenden gemacht werden, denn nur unter Voraussetzung einer ganz gleichen Behandlung ist es möglich geworden, das Deficit auf den mässigen Rest von 2963 Rthl. herabzubringen. Zudem würden die neuen Witwen, wenn den frühern eine exceptionelle Bevorzugung \textit{aus der Kasse} gewährt würde, sich für lädirt halten, da die Möglichkeit, bei so sehr gelinden Modificationen der Pensionen stehen zu bleiben, bloss durch die Erhöhung der Beiträge bewirkt wird, an welcher die Ehemänner der neuen Witwen Theil genommen haben werden, die der frühern aber nicht. Da jedoch, von der andern Seite, als wunschenswerth erscheint, dass jeder wenn auch im Rechte nicht begründeten Unzufriedenheit der gegenwärtigen Witwen vorgebeugt werde. so schlägt die Commission vor
\begin{quote}
dass von Seiten der Universität an das h. Curatorium die Bitte gerichtet werde, für die Zeit, wo in Folge obiger Regulirung der Pensionsbetrag unter 250 Rthl. herabgehen wird, den gegenwärtigen Witwen, so viele dann noch am Leben sein werden, das an 250 Rthl. fehlende aus der Universitätskasse ergänzen zu lassen.
\end{quote}
\noindent\hspace*{\parindent}Eine solche Bitte wurde sich mit nahe liegenden Gründen unterstutzen lassen. Auch ist leicht zu übersehen, dass eine solche Beihülfe der Universitätskasse keine grosse Last auflegen konne. Unmöglich wäre es sogar nicht, dass der Fall der Beihülfe \textit{niemals} einträte. Aber sollte sie auch, wenn früher eine neue Witwe hinzukommt, ehe eine der bisherigen abgegangen ist, bald schon erforderlich werden, so wird doch voraussichtlich für geraume Zeit der Zuschuss fur jede einzelne Witwe nur wenige Thaler betragen konnen. In späterer Zukunft, z. B. nach 30 Jahren, konnte wenn wir Beispielshalber einen der extremsten Fälle annehmen, dass nemlich bis dahin gar keine Erhöhung der Pension zulässig geworden und die Zahl der Witwen auf 30 gestiegen wäre, die Grösse des Zuschusses fur eine Witwe 20 Rthl. betragen: allein dann werden wahrscheinlich von den jetzigen Witwen nur noch wenige am Leben sein.

Die Commission glaubt hier noch ein paar andere Einrichtungen, die als Mittel zur Verbesserung der Bilanz zur Sprache gebracht sind, erwähnen zu müssen, ohne jedoch einen Antrag darauf richten zu wollen.

Es ist zuvörderst in Frage gekommen, ob nicht einige Verbesserung der Bilanz dadurch erreicht werden konne, dass den künftig eintretenden ausserordentlichen Professoren nur der Anspruch auf eine geringere Witwenpension beigelegt wurde. Diese würden sonach eine zweite Klasse von Interessenten bilden, aus der sie bei ihrer Beförderung zur ordentlichen Professur von selbst in die erste hinaufrücken: die Beiträge der Mitglieder zweiter Klasse sollten dagegen ihre bisherige Grösse von 2 Louisd’or behalten. Den jetzigen ausserordentlichen Professoren sollte die Wahl gelassen werden, ob sie in der ersten Klasse bleiben oder in die zweite übertreten wollten, so jedoch, dass im letztern Fall ein Rücktritt in die erste Klasse nicht zulässig wäre, so lange sie ausserordentliche Professoren blieben.

Zu richtiger Würdigung dieses Gedankens ist zuvörderst wohl zu bedenken, dass in dem oben S. [160] aufgestellten Hauptplan die Reduction des Deficits von 17520 Rthl. auf 2963 Rthl. wesentlich von der Voraussetzung abhängt, dass die Beiträge \textit{aller} Mitglieder, der jetzigen wie der künftigen von 2 auf 3 Louisd’or erhöhet werden, und dass mithin jener Hauptplan gar nicht bestehen kann, wenn diese Voraussetzung einen wesentlichen Abgang erleidet, ohne dass dafür anderweit ein vollständiger und sicherer Ersatz eintritt. Durch diese Betrachtung wird eine sonst durch ihre Einfachheit sich empfehlende Art, die Pension der zweiten Klasse niedriger zu normiren, von selbst als unzulässig ausgeschlossen, die nemlich, dass man den Witwen zweiter Klasse nur den fixen Theil der Pension, nemlich 200 Rthl. einraume, oder sie völlig der Professorin G. gleichstellen sollte. Denn in der That ist leicht zu übersehen, dass dadurch sehr wahrscheinlich \textit{die Kasse} gar keinen Ersatz für jene Einbusse an Beiträgen erhalten, und nur die Witwen erster Klasse etwas besser gestellt werden, indem die Zusatzdividende unter eine geringere Zahl von Participanten vertheilt wurde.
